\clearpage
\chapter*{Introduction générale}
\markboth{Introduction générale}{Introduction générale}
\addtocontents{toc}{\protect\contentsline{introduction}{Introduction générale}{\thepage}{intro.1}}

Dans un contexte professionnel en constante mutation, où la transformation numérique et l'intelligence artificielle redéfinissent les modes d'apprentissage et de gestion des compétences, la formation continue devient un levier stratégique essentiel pour les entreprises. L'émergence de l'intelligence artificielle générative et des systèmes multi-agents ouvre aujourd'hui de nouvelles perspectives pour concevoir des plateformes d'apprentissage plus intelligentes, adaptatives et personnalisées.

C'est dans cette dynamique que s'inscrit notre projet de fin d'études, réalisé au sein de la société Standard Sharing Software (3S), acteur majeur de l'intégration de solutions technologiques en Tunisie et en Afrique. Ce projet vise à concevoir et développer une plateforme web intelligente de formation et de mentorat, capable de centraliser la gestion des compétences des collaborateurs du Groupe 3S tout en offrant un accompagnement pédagogique personnalisé basé sur l'intelligence artificielle.

La solution proposée repose sur une architecture microservices multi-tenant intégrant plusieurs agents intelligents spécialisés, chacun ayant un rôle précis dans le processus d'apprentissage : génération de contenus pédagogiques, coaching personnalisé, recommandation de ressources, suivi des apprenants et gestion de la mémoire contextuelle. Ces agents s'appuient sur des modèles de langage avancés, un système de Retrieval-Augmented Generation (RAG) agentique, ainsi qu'un mécanisme de mémoire persistante permettant d'assurer un suivi cohérent et évolutif des utilisateurs.

Ce rapport présente l'ensemble des travaux réalisés au cours de ce projet et s'articule autour de cinq chapitres principaux, chacun mettant en évidence une étape clé de la démarche.

Le premier chapitre, intitulé « Contexte général et méthodologie de travail », présente le cadre du projet, l'organisme d'accueil ainsi que la problématique étudiée. Il expose également les exigences fonctionnelles et non fonctionnelles, ainsi que la méthodologie de développement adoptée.

Le deuxième chapitre, « Concepts de base et étude préalable », introduit les fondements théoriques liés à l'intelligence artificielle générative, aux grands modèles de langage, aux systèmes multi-agents et aux approches de type RAG, qui constituent la base scientifique du projet.

Le troisième chapitre, « Conception et architecture de la plateforme », décrit l'architecture globale du système, notamment l'organisation en microservices, le modèle multi-tenant, ainsi que la conception des agents intelligents et les choix technologiques retenus.

Le quatrième chapitre, « Réalisation », présente les différents modules développés ainsi que les fonctionnalités implémentées, en mettant l'accent sur les composants liés à l'intelligence artificielle, à la gestion des formations et à l'expérience utilisateur.

Le cinquième chapitre, « Tests et validation », détaille la stratégie de test adoptée ainsi que les résultats obtenus à travers les tests unitaires, fonctionnels et de performance, ainsi que l'évaluation des modules d'intelligence artificielle.

Enfin, la conclusion générale synthétisera les apports du projet, évaluera le niveau d'atteinte des objectifs fixés et proposera des perspectives d'amélioration et d'évolution future de la plateforme.
